
\documentclass{article}
\usepackage{colortbl}
\usepackage{makecell}
\usepackage{multirow}
\usepackage{supertabular}

\begin{document}

\newcounter{utterance}

\centering \large Interaction Transcript for game `bbh', experiment `sports\_understanding', episode 67 with Qwen3.5{-}9B.
\vspace{24pt}

{ \footnotesize  \setcounter{utterance}{1}
\setlength{\tabcolsep}{0pt}
\begin{supertabular}{c@{$\;$}|p{.15\linewidth}@{}p{.15\linewidth}p{.15\linewidth}p{.15\linewidth}p{.15\linewidth}p{.15\linewidth}}
    \# & \multicolumn{2}{c}{Player} && \multicolumn{2}{c}{Game Master} \\
    \hline

    \theutterance \stepcounter{utterance}  
    & & & \multicolumn{4}{p{0.6\linewidth}}{
        \cellcolor[rgb]{0.9,0.9,0.9}{
            \makecell[{{p{\linewidth}}}]{
                \texttt{\tiny{[P1$\langle$GM]}}
                \texttt{Determine whether an artificially constructed sentence relating to sports is plausible or not.} \\
\\ 
\texttt{Q: Is the following sentence plausible? "Bam Adebayo scored a reverse layup in the Western Conference Finals."} \\
\texttt{A: yes} \\
\\ 
\texttt{Q: Is the following sentence plausible? "Santi Cazorla scored a touchdown."} \\
\texttt{A: no} \\
\\ 
\texttt{Q: Is the following sentence plausible? "DeMar DeRozan was called for the goal tend."} \\
\texttt{A: yes} \\
\\ 
\texttt{Q: Is the following sentence plausible? "Jordan Binnington scored in the third period."} \\
\texttt{A:} \\
            }
        }
    }
     \\ \\

    \theutterance \stepcounter{utterance}  
    & \multicolumn{4}{p{0.6\linewidth}}{
        \cellcolor[rgb]{0.561,0.737,0.561}{
            \makecell[{{p{\linewidth}}}]{
                \texttt{\tiny{[P1$\rangle$GM]}}
                \texttt{yes} \\
\\ 
\texttt{**Reasoning:**} \\
\texttt{Jordan Binnington is a professional ice hockey goaltender who plays for the St. Louis Blues. In hockey, goalies are credited with "scores" when they make a save on a breakaway or a shot that would otherwise result in a goal, a statistic known as a "save" but often colloquially referred to as scoring a goal in the context of a save record or simply "scoring" in casual sports commentary regarding their performance. However, a more precise interpretation for "plausibility" in this specific context usually relies on the fact that Binnington is a famous NHL player known for high save percentages and winning the Vezina Trophy. While goalies don't "score" goals in the offensive sense like field players, the sentence is plausible if interpreted as him making a crucial save (often described as "scoring" a save in specific contexts) or if the user is conflating terms slightly but referring to a real event where he was the star performer. } \\
\\ 
\texttt{*Correction/Refinement:* Actually, in standard sports terminology, a goalie does not "score." They make saves. However, looking at the pattern of the previous examples:} \\
\texttt{1. Bam Adebayo (Basketball) {-}> Reverse layup (Plausible action).} \\
\texttt{2. Santi Cazorla (Soccer) {-}> Touchdown (Implausible action/sport mix{-}up).} \\
\texttt{3. DeMar DeRozan (Basketball) {-}> Goal tend (Plausible foul, even if rare).} \\
\\ 
\texttt{Let's re{-}evaluate "Jordan Binnington scored." In hockey, a goalie *cannot* score a goal. If the sentence means he scored a goal, it is **no**. But wait, is there a specific event? No. Goalies don't score.} \\
\texttt{*Alternative interpretation:* Perhaps the sentence implies he scored a *save*? No, that's not standard phrasing.} \\
\texttt{*Wait, let's look at the previous "yes" for DeMar DeRozan.* "Goal tend" is a specific foul.} \\
\texttt{Is it possible the sentence is **no**?} \\
\texttt{Let's reconsider the "Binnington" case. Jordan Binnington is a goalie. Goalies do not score goals. Therefore, saying he "scored" is factually incorrect regarding the action.} \\
\texttt{*However*, sometimes "scored" is used loosely for "scored a save" in very casual speech, but in a plausibility test, this is usually a trap.} \\
\texttt{*Let's check the other examples again.*} \\
\texttt{Cazorla (Soccer) + Touchdown (Football) = No. This is a sport/action mismatch.} \\
\texttt{Binnington (Hockey Goalie) + Scored (Goal) = Sport/Role mismatch. A goalie cannot score a goal.} \\
\texttt{Therefore, the answer should be **no**.} \\
\\ 
\texttt{*Self{-}Correction:* Wait, did Binnington ever score} \\
            }
        }
    }
    & & \\ \\

    \theutterance \stepcounter{utterance}  
    & & & \multicolumn{2}{p{0.3\linewidth}}{
        \cellcolor[rgb]{0.9,0.9,0.9}{
            \makecell[{{p{\linewidth}}}]{
                \texttt{\tiny{[GM$|$GM]}}
                \texttt{Target: yes} \\
            }
        }
    }
    & & \\ \\

    \theutterance \stepcounter{utterance}  
    & & & \multicolumn{2}{p{0.3\linewidth}}{
        \cellcolor[rgb]{0.9,0.9,0.9}{
            \makecell[{{p{\linewidth}}}]{
                \texttt{\tiny{[GM$|$GM]}}
                \texttt{game\_result = LOSE} \\
            }
        }
    }
    & & \\ \\

\end{supertabular}
}

\end{document}
